% assets/a3-pure-sway.svg — 순수 횡이동의 닫힌 해와 요 모멘트가 0 이 되는 이유
%
%   bash _tools/tikz2svg.sh assets/src/a3-pure-sway.tex assets/a3-pure-sway.svg
%
% 이 그림이 말하는 것
%   tau = [0; Y; 0] 을 넣으면 pseudo-inverse 해가 한 점으로 떨어진다.
%       u_e = (Y/2) [L/b, 1, -L/b, 1]'
%   이때 두 추력선이 **둘 다 무게중심을 지난다**. 모멘트 팔이 0 이므로
%   요 모멘트는 상쇄되는 것이 아니라 **애초에 생기지 않는다**.
%       tan(delta) = b/L  <=>  (-L, -b) 에서 (L, b) 방향 = 원점을 지나는 직선
%   실측에서 선수각 RMS 가 0.0000 도로 나온 것이 이 기하 때문이다.
%
% 좌표 — a3-geometry.tex 와 같은 규약. 1 m = 20 mm
%   화면 x = y_b (우현), 화면 y = x_b (선수)
%
% 우현 추진기는 **역추력**이다 (T_R < 0)
%   노즐은 -23.4 도(선수-좌현 쪽)를 보고 있는데 추력이 음수라, 실제로 나는 힘은
%   그 반대인 선미-우현 쪽이다. 노즐 방향은 가는 선, 실제 힘은 굵은 화살표로 나눠 그린다
%
% 겹침을 피하려고 정한 것
%   좌현 화살표 1.8 m, 우현 화살표 1.2 m — 길이를 다르게 해 라벨이 겹치지 않게 한다
%   무게중심을 지나는 두 작용선은 점선으로 선체 밖까지 늘여 그린다

\documentclass[border=6pt]{standalone}
% capstone-style.tex — 이 볼트의 TikZ 그림이 공유하는 스타일
%
% 쓰는 법:  % capstone-style.tex — 이 볼트의 TikZ 그림이 공유하는 스타일
%
% 쓰는 법:  % capstone-style.tex — 이 볼트의 TikZ 그림이 공유하는 스타일
%
% 쓰는 법:  \input{capstone-style}   (assets/src/ 안에서)
% 빌드:     bash _tools/tikz2svg.sh assets/src/w02-node-topic.tex assets/w02-node-topic.svg
%
% 왜 TikZ 인가
%   손으로 SVG 좌표를 쓰면 화살촉·선 굵기·글자 크기가 그림마다 미묘하게 달라진다.
%   그것이 "자료가 어설퍼 보인다"의 정체다. 스타일을 한 번만 정하고 상대좌표로
%   배치하면 좌표를 고쳐도 그림이 무너지지 않는다.
%   수식은 본문과 같은 TeX 글꼴로 나온다.
%
% 한글
%   ko.TeX(DVI 경로)를 쓴다. XeLaTeX 은 TikZ 그래픽을 PDF special 로 내보내는데
%   dvisvgm 이 그것을 받으려면 Ghostscript < 10.01 이 필요하다. 이 PC 에는
%   10.07 이 깔려 있어 그 경로가 막힌다. latex -> DVI -> dvisvgm 이 유일하게
%   동작하는 조합이다.

\usepackage[hangul]{kotex}
\usepackage{tikz}
\usepackage{amsmath}
\usepackage{xcolor}
\usetikzlibrary{arrows.meta, positioning, calc, fit, backgrounds, decorations.pathmorphing}

% ---- 색 --------------------------------------------------------------------
% 강조는 하나만. 나머지는 검정. 흑백 인쇄에서도 읽혀야 한다.
\definecolor{ink}{HTML}{1A1A1A}
\definecolor{soft}{HTML}{666666}
\definecolor{accent}{HTML}{7C3AED}
\definecolor{warn}{HTML}{B91C1C}

% 계층 색 — Simulink 모델의 블록 색과 같은 규칙을 쓴다
\definecolor{guid}{HTML}{1D4ED8}    % 유도
\definecolor{ctrl}{HTML}{EA580C}    % 제어
\definecolor{thr}{HTML}{D97706}     % 추진기
\definecolor{plant}{HTML}{16A34A}   % 운동모델
\definecolor{nav}{HTML}{7C3AED}     % 항법 · 신호처리
\definecolor{pathc}{HTML}{0D9488}   % 계획 경로 · 항적
\definecolor{errc}{HTML}{C2410C}    % 오차
\definecolor{flow}{HTML}{0369A1}    % 조류 · 바람

% ---- 선과 화살촉 -----------------------------------------------------------
\tikzset{
  sig/.style   = {ink, line width=0.5pt, -{Stealth[length=4pt, width=3pt]}},
  wire/.style  = {ink, line width=0.5pt},
  fb/.style    = {ink, line width=0.5pt, densely dashed,
                  -{Stealth[length=4pt, width=3pt]}},
  asig/.style  = {accent, line width=0.5pt, densely dashed,
                  -{Stealth[length=4pt, width=3pt]}},
  % 블록 — 흰 바탕, 직각, 검은 테두리. 색은 테두리에만 준다
  blk/.style   = {draw=ink, line width=0.55pt, fill=white,
                  minimum width=18mm, minimum height=8.5mm, inner sep=2pt, font=\small},
  gblk/.style  = {blk, draw=guid},
  cblk/.style  = {blk, draw=ctrl},
  tblk/.style  = {blk, draw=thr},
  pblk/.style  = {blk, draw=plant},
  nblk/.style  = {blk, draw=nav},
  % 합산점 · 분기점
  sum/.style   = {draw=ink, line width=0.55pt, fill=white, circle,
                  inner sep=0pt, minimum size=4.6mm, font=\scriptsize},
  dot/.style   = {ink, fill=ink, circle, inner sep=0pt, minimum size=1.5mm},
  % 글자
  sname/.style = {font=\small, inner sep=1.5pt},
  note/.style  = {font=\scriptsize, soft, inner sep=1.5pt},
  wnote/.style = {font=\scriptsize, warn, inner sep=1.5pt},
  anote/.style = {font=\scriptsize, accent, inner sep=1.5pt},
  axis/.style  = {ink, line width=0.5pt},
  tick/.style  = {font=\scriptsize, soft},
  % 묶음 상자 — 서브시스템 표시
  grp/.style   = {draw=soft, line width=0.4pt, densely dotted, rounded corners=1mm,
                  inner sep=3mm},
}
   (assets/src/ 안에서)
% 빌드:     bash _tools/tikz2svg.sh assets/src/w02-node-topic.tex assets/w02-node-topic.svg
%
% 왜 TikZ 인가
%   손으로 SVG 좌표를 쓰면 화살촉·선 굵기·글자 크기가 그림마다 미묘하게 달라진다.
%   그것이 "자료가 어설퍼 보인다"의 정체다. 스타일을 한 번만 정하고 상대좌표로
%   배치하면 좌표를 고쳐도 그림이 무너지지 않는다.
%   수식은 본문과 같은 TeX 글꼴로 나온다.
%
% 한글
%   ko.TeX(DVI 경로)를 쓴다. XeLaTeX 은 TikZ 그래픽을 PDF special 로 내보내는데
%   dvisvgm 이 그것을 받으려면 Ghostscript < 10.01 이 필요하다. 이 PC 에는
%   10.07 이 깔려 있어 그 경로가 막힌다. latex -> DVI -> dvisvgm 이 유일하게
%   동작하는 조합이다.

\usepackage[hangul]{kotex}
\usepackage{tikz}
\usepackage{amsmath}
\usepackage{xcolor}
\usetikzlibrary{arrows.meta, positioning, calc, fit, backgrounds, decorations.pathmorphing}

% ---- 색 --------------------------------------------------------------------
% 강조는 하나만. 나머지는 검정. 흑백 인쇄에서도 읽혀야 한다.
\definecolor{ink}{HTML}{1A1A1A}
\definecolor{soft}{HTML}{666666}
\definecolor{accent}{HTML}{7C3AED}
\definecolor{warn}{HTML}{B91C1C}

% 계층 색 — Simulink 모델의 블록 색과 같은 규칙을 쓴다
\definecolor{guid}{HTML}{1D4ED8}    % 유도
\definecolor{ctrl}{HTML}{EA580C}    % 제어
\definecolor{thr}{HTML}{D97706}     % 추진기
\definecolor{plant}{HTML}{16A34A}   % 운동모델
\definecolor{nav}{HTML}{7C3AED}     % 항법 · 신호처리
\definecolor{pathc}{HTML}{0D9488}   % 계획 경로 · 항적
\definecolor{errc}{HTML}{C2410C}    % 오차
\definecolor{flow}{HTML}{0369A1}    % 조류 · 바람

% ---- 선과 화살촉 -----------------------------------------------------------
\tikzset{
  sig/.style   = {ink, line width=0.5pt, -{Stealth[length=4pt, width=3pt]}},
  wire/.style  = {ink, line width=0.5pt},
  fb/.style    = {ink, line width=0.5pt, densely dashed,
                  -{Stealth[length=4pt, width=3pt]}},
  asig/.style  = {accent, line width=0.5pt, densely dashed,
                  -{Stealth[length=4pt, width=3pt]}},
  % 블록 — 흰 바탕, 직각, 검은 테두리. 색은 테두리에만 준다
  blk/.style   = {draw=ink, line width=0.55pt, fill=white,
                  minimum width=18mm, minimum height=8.5mm, inner sep=2pt, font=\small},
  gblk/.style  = {blk, draw=guid},
  cblk/.style  = {blk, draw=ctrl},
  tblk/.style  = {blk, draw=thr},
  pblk/.style  = {blk, draw=plant},
  nblk/.style  = {blk, draw=nav},
  % 합산점 · 분기점
  sum/.style   = {draw=ink, line width=0.55pt, fill=white, circle,
                  inner sep=0pt, minimum size=4.6mm, font=\scriptsize},
  dot/.style   = {ink, fill=ink, circle, inner sep=0pt, minimum size=1.5mm},
  % 글자
  sname/.style = {font=\small, inner sep=1.5pt},
  note/.style  = {font=\scriptsize, soft, inner sep=1.5pt},
  wnote/.style = {font=\scriptsize, warn, inner sep=1.5pt},
  anote/.style = {font=\scriptsize, accent, inner sep=1.5pt},
  axis/.style  = {ink, line width=0.5pt},
  tick/.style  = {font=\scriptsize, soft},
  % 묶음 상자 — 서브시스템 표시
  grp/.style   = {draw=soft, line width=0.4pt, densely dotted, rounded corners=1mm,
                  inner sep=3mm},
}
   (assets/src/ 안에서)
% 빌드:     bash _tools/tikz2svg.sh assets/src/w02-node-topic.tex assets/w02-node-topic.svg
%
% 왜 TikZ 인가
%   손으로 SVG 좌표를 쓰면 화살촉·선 굵기·글자 크기가 그림마다 미묘하게 달라진다.
%   그것이 "자료가 어설퍼 보인다"의 정체다. 스타일을 한 번만 정하고 상대좌표로
%   배치하면 좌표를 고쳐도 그림이 무너지지 않는다.
%   수식은 본문과 같은 TeX 글꼴로 나온다.
%
% 한글
%   ko.TeX(DVI 경로)를 쓴다. XeLaTeX 은 TikZ 그래픽을 PDF special 로 내보내는데
%   dvisvgm 이 그것을 받으려면 Ghostscript < 10.01 이 필요하다. 이 PC 에는
%   10.07 이 깔려 있어 그 경로가 막힌다. latex -> DVI -> dvisvgm 이 유일하게
%   동작하는 조합이다.

\usepackage[hangul]{kotex}
\usepackage{tikz}
\usepackage{amsmath}
\usepackage{xcolor}
\usetikzlibrary{arrows.meta, positioning, calc, fit, backgrounds, decorations.pathmorphing}

% ---- 색 --------------------------------------------------------------------
% 강조는 하나만. 나머지는 검정. 흑백 인쇄에서도 읽혀야 한다.
\definecolor{ink}{HTML}{1A1A1A}
\definecolor{soft}{HTML}{666666}
\definecolor{accent}{HTML}{7C3AED}
\definecolor{warn}{HTML}{B91C1C}

% 계층 색 — Simulink 모델의 블록 색과 같은 규칙을 쓴다
\definecolor{guid}{HTML}{1D4ED8}    % 유도
\definecolor{ctrl}{HTML}{EA580C}    % 제어
\definecolor{thr}{HTML}{D97706}     % 추진기
\definecolor{plant}{HTML}{16A34A}   % 운동모델
\definecolor{nav}{HTML}{7C3AED}     % 항법 · 신호처리
\definecolor{pathc}{HTML}{0D9488}   % 계획 경로 · 항적
\definecolor{errc}{HTML}{C2410C}    % 오차
\definecolor{flow}{HTML}{0369A1}    % 조류 · 바람

% ---- 선과 화살촉 -----------------------------------------------------------
\tikzset{
  sig/.style   = {ink, line width=0.5pt, -{Stealth[length=4pt, width=3pt]}},
  wire/.style  = {ink, line width=0.5pt},
  fb/.style    = {ink, line width=0.5pt, densely dashed,
                  -{Stealth[length=4pt, width=3pt]}},
  asig/.style  = {accent, line width=0.5pt, densely dashed,
                  -{Stealth[length=4pt, width=3pt]}},
  % 블록 — 흰 바탕, 직각, 검은 테두리. 색은 테두리에만 준다
  blk/.style   = {draw=ink, line width=0.55pt, fill=white,
                  minimum width=18mm, minimum height=8.5mm, inner sep=2pt, font=\small},
  gblk/.style  = {blk, draw=guid},
  cblk/.style  = {blk, draw=ctrl},
  tblk/.style  = {blk, draw=thr},
  pblk/.style  = {blk, draw=plant},
  nblk/.style  = {blk, draw=nav},
  % 합산점 · 분기점
  sum/.style   = {draw=ink, line width=0.55pt, fill=white, circle,
                  inner sep=0pt, minimum size=4.6mm, font=\scriptsize},
  dot/.style   = {ink, fill=ink, circle, inner sep=0pt, minimum size=1.5mm},
  % 글자
  sname/.style = {font=\small, inner sep=1.5pt},
  note/.style  = {font=\scriptsize, soft, inner sep=1.5pt},
  wnote/.style = {font=\scriptsize, warn, inner sep=1.5pt},
  anote/.style = {font=\scriptsize, accent, inner sep=1.5pt},
  axis/.style  = {ink, line width=0.5pt},
  tick/.style  = {font=\scriptsize, soft},
  % 묶음 상자 — 서브시스템 표시
  grp/.style   = {draw=soft, line width=0.4pt, densely dotted, rounded corners=1mm,
                  inner sep=3mm},
}


\begin{document}
\begin{tikzpicture}[x=1mm, y=1mm]

\pgfmathsetmacro{\BS}{20}
\newcommand{\BQ}[2]{{(#2)*\BS},{(#1)*\BS}}
\newcommand{\BP}[2]{(\BQ{#1}{#2})}

\pgfmathsetmacro{\bb}{1.027135}
\pgfmathsetmacro{\LL}{2.373776}

% ===========================================================================
% (가) 두 작용선이 무게중심에서 만난다
% ===========================================================================
\begin{scope}[shift={(0,0)}]

  \node[font=\small\bfseries, anchor=north west] at (-40,62)
       {(가) 두 추력선이 모두 무게중심을 지난다};

  % --- 선체 ---------------------------------------------------------------
  \draw[ink, line width=0.6pt, fill=white]
       \BP{-2.75}{-\bb} -- \BP{1.3}{-\bb} -- \BP{2.4}{0} -- \BP{1.3}{\bb}
       -- \BP{-2.75}{\bb} -- cycle;
  \draw[soft, line width=0.3pt, densely dashed] \BP{-2.85}{0} -- \BP{2.5}{0};
  \node[font=\scriptsize, nav, anchor=south] at \BP{2.45}{0} {$x_b$};

  % --- 작용선 (무게중심을 지난다) ----------------------------------------
  \draw[plant, line width=0.5pt, densely dotted] \BP{-2.70}{-1.168} -- \BP{0.80}{0.346};
  \draw[plant, line width=0.5pt, densely dotted] \BP{-3.60}{1.557}  -- \BP{0.80}{-0.346};

  % --- 좌현 추진기 — 노즐 +23.4 도, 추력 양수, 길이 1.8 m ----------------
  %     끝점 x_b = -0.7220,  y_b = -0.3119
  \draw[thr, line width=1.4pt, -{Stealth[length=4.6pt,width=3.4pt]}]
       \BP{-\LL}{-\bb} -- \BP{-0.7220}{-0.3119};
  \node[font=\scriptsize, thr, anchor=east] at \BP{-1.30}{-0.80} {$T_L > 0$};
  \node[circle, draw=thr, fill=white, line width=0.7pt, minimum size=3.4mm,
        inner sep=0pt] at \BP{-\LL}{-\bb} {};
  \node[font=\scriptsize, ctrl, anchor=north east] at \BP{-2.52}{-1.10}
       {$\delta_L = +23.40^\circ$};

  % --- 우현 추진기 — 노즐 -23.4 도, 추력 음수(역추력), 길이 1.2 m --------
  %     노즐이 보는 쪽 (가는 선):  x_b -1.6727, y_b +0.6500  (0.75 m)
  %     실제 힘   (굵은 화살표):  x_b -3.4750, y_b +1.5039  (1.2 m, 반대쪽)
  \draw[soft, line width=0.8pt] \BP{-\LL}{\bb} -- \BP{-1.6727}{0.6500};
  \node[font=\scriptsize, soft, anchor=south west] at \BP{-1.70}{0.64} {노즐이 보는 쪽};
  \draw[thr, line width=1.4pt, -{Stealth[length=4.6pt,width=3.4pt]}]
       \BP{-\LL}{\bb} -- \BP{-3.4750}{1.5039};
  \node[font=\scriptsize, thr, anchor=west] at \BP{-3.20}{1.55} {$T_R < 0$ (역추력)};
  \node[circle, draw=thr, fill=white, line width=0.7pt, minimum size=3.4mm,
        inner sep=0pt] at \BP{-\LL}{\bb} {};
  \node[font=\scriptsize, ctrl, anchor=north west] at \BP{-2.52}{1.10}
       {$\delta_R = -23.40^\circ$};

  % --- 무게중심과 합력 ----------------------------------------------------
  \draw[plant, line width=1.6pt, -{Stealth[length=5pt,width=3.6pt]}]
       \BP{0}{0} -- \BP{0}{1.45};
  \node[font=\scriptsize, plant, anchor=south] at \BP{0.08}{0.80} {합력 $Y$};
  \node[dot] at \BP{0}{0} {};
  \node[font=\scriptsize, soft, anchor=east] at \BP{-0.14}{-0.10} {무게중심};

  % --- 모멘트 팔이 0 이라는 표시 -----------------------------------------
  \node[anchor=north west, font=\scriptsize, plant, text width=86mm] at (-40,-78)
       {\textbf{모멘트 팔이 $0$ 이다} — 상쇄되는 것이 아니라 생기지 않는다 \\[3pt]
        $\tan\delta = b/L$ 이라는 것은, $(-L,\,-b)$ 에 앉은 추진기가
        $(L,\,b)$ 방향을 본다는 뜻이다. 그 직선은 반드시 원점을 지난다};

  \node[anchor=north west, font=\scriptsize, text width=86mm] at (-40,-94)
       {\textbf{읽는 법} \\[3pt]
        좌현은 앞으로, 우현은 뒤로 민다 — 앞뒤 몫은 서로 지워져 $X = 0$ \\[2pt]
        옆으로 미는 몫은 둘 다 우현 쪽이라 더해져 $Y$ 가 된다 \\[2pt]
        두 작용선이 무게중심에서 만나므로 $N = 0$};
\end{scope}

% ===========================================================================
% (나) 닫힌 해 — 행렬 한 줄에서 나온다
% ===========================================================================
\begin{scope}[shift={(72,0)}]

  \node[font=\small\bfseries, anchor=north west] at (0,62)
       {(나) 요구 $\boldsymbol\tau = [0,\ Y,\ 0]^{\top}$ 을 넣으면 해가 한 점으로 떨어진다};

  \node[anchor=north, font=\small] at (52,52)
       {$\mathbf{u}_e = \mathbf{B}^{+}\begin{bmatrix} 0 \\ Y \\ 0 \end{bmatrix}
         = \frac{Y}{2}\begin{bmatrix}
             \textcolor{guid}{L/b} \\ \textcolor{guid}{1} \\
             \textcolor{ctrl}{-L/b} \\ \textcolor{ctrl}{1}
           \end{bmatrix}
         = Y\begin{bmatrix}
             \textcolor{guid}{+1.1554} \\ \textcolor{guid}{+0.5} \\
             \textcolor{ctrl}{-1.1554} \\ \textcolor{ctrl}{+0.5}
           \end{bmatrix}$};

  \draw[soft, line width=0.3pt] (0,24) -- (104,24);

  \node[anchor=north west, font=\scriptsize, text width=104mm] at (0,19)
       {\textbf{극좌표로 되돌리면 각도와 추력이 수치로 나온다} \\[4pt]
        \hspace*{6mm}$\delta_L = \operatorname{atan2}\!\left(\tfrac{Y}{2},\ \tfrac{Y}{2}\tfrac{L}{b}\right)
        = \arctan\dfrac{b}{L} = +23.3982^\circ$,
        \qquad $\delta_R = -23.3982^\circ$ \\[5pt]
        \hspace*{6mm}$T_L = -T_R = \dfrac{\sqrt{b^2 + L^2}}{2b}\,Y = 1.259069\,Y$ \\[5pt]
        $Y$ 하나만 정하면 각도 둘과 추력 둘이 \textbf{전부 결정된다} — 조절할 것이 없다};

  \draw[soft, line width=0.3pt] (0,-16) -- (104,-16);

  \node[anchor=north west, font=\scriptsize, text width=104mm] at (0,-21)
       {\textbf{요 모멘트가 0 인 것을 식으로 확인한다} \\[4pt]
        \hspace*{2mm}$N = \underbrace{b F_{x,L} - L F_{y,L}}_{%
            \textstyle +\frac{LY}{2} - \frac{LY}{2} = 0}
            \;+\; \underbrace{-\,b F_{x,R} - L F_{y,R}}_{%
            \textstyle +\frac{LY}{2} - \frac{LY}{2} = 0} \;=\; 0$ \\[6pt]
        좌우가 서로 지우는 것이 아니라 \textbf{각자가 이미 0} 이다.
        그래서 좌우 추진기 성능이 조금 달라도 선수각이 흔들리지 않는다};

  \draw[soft, line width=0.3pt] (0,-56) -- (104,-56);

  \node[anchor=north west, font=\scriptsize, text width=104mm] at (0,-61)
       {\textbf{추력 한계가 횡이동 속도의 천장을 정한다} \\[4pt]
        \hspace*{6mm}$\lvert T \rvert \le F_{\max} = 500$ N
        \;$\Longrightarrow$\;
        $Y_{\max} = \dfrac{500}{1.259069} = 397.1$ N \\[4pt]
        기준 환경 실측 — 정상 횡이동 속도 $1.555$ m/s (닫힌식과 차이 $0.000$)};

  \node[anchor=north west, font=\scriptsize, warn, text width=104mm] at (0,-84)
       {\textbf{순수 횡이동은 위치 제어의 한 경우다.} 배분기도 운동모델도 그대로이고,
        설정값만 $x_{\mathrm{ref}}$ 고정 $\cdot$ $y_{\mathrm{ref}}$ 이동 $\cdot$
        $\psi_{\mathrm{ref}}$ 고정으로 둔 것뿐이다 — 코드는 한 줄도 다르지 않다};
\end{scope}

\end{tikzpicture}
\end{document}
